%
\documentclass[final,notitlepage,12pt,reqno]{amsart}

%\usepackage{arydshln}
\usepackage[table]{xcolor}
\usepackage{etoolbox,mathtools}
\makeatletter
\patchcmd{\@makefnmark}{\fontsize}{\check@mathfonts\fontsize}{}{}
\makeatother

\makeatletter

\DeclareFontFamily{U}{mathx}{}
\DeclareFontShape{U}{mathx}{m}{n}{<-> mathx10}{}
\DeclareSymbolFont{mathx}{U}{mathx}{m}{n}
\DeclareMathAccent{\widecheck}{0}{mathx}{"71}

\def\smallunderbrace#1{\mathop{\vtop{\m@th\ialign{##\crcr
   $\hfil\displaystyle{#1}\hfil$\crcr
   \noalign{\kern3\p@\nointerlineskip}%
   \tiny\upbracefill\crcr\noalign{\kern3\p@}}}}\limits}
\makeatother

\newcommand{\bmmax}{0}
\newcommand{\hmmax}{0}
\UseRawInputEncoding
\usepackage{graphicx,extarrows}
\usepackage{calc}
\usepackage{url}
\newlength{\depthofsumsign}
\setlength{\depthofsumsign}{\depthof{$\sum$}}
\newlength{\totalheightofsumsign}
\newlength{\heightanddepthofargument}

\newcommand{\nsum}[1][1.0]{% only for \displaystyle
   \mathop{%
        \raisebox
            {-#1\depthofsumsign+1\depthofsumsign}
            {\scalebox
                {#1}
                {$\displaystyle\sum$}%
            }
    }
}



\makeatletter
\let\I\@undefined
\makeatother

\setlength{\emergencystretch}{0em}

\usepackage{geometry}
\geometry{body={17.5cm,24.5cm}}


\usepackage{indentfirst}
\usepackage[normalem]{ulem}
\usepackage{float}
\usepackage{amsthm}

\usepackage{enumitem}[2011/09/28]
\setenumerate{align=left, leftmargin=0pt,labelsep=.5em, labelindent=0\parindent,listparindent=\parindent,itemindent=*}
%\setitemize{noitemsep,topsep=0pt,parsep=0pt,partopsep=0pt}
\usepackage{array}
%\usepackage{empheq}
%\usepackage{natbib}
%\setlength{\bibsep}{0.0pt}
\usepackage[all]{xy}


\usepackage{exscale,relsize}
\usepackage{multirow}


\usepackage{caption}[2012/02/19]

\usepackage{longtable,lscape}

%\captionsetup[figure]{font=small}
\captionsetup[table]{font=small}
\setlength{\LTcapwidth}{\textwidth}
\usepackage{colonequals}

\usepackage{amssymb,bm,amsmath}
\usepackage{amsfonts}

\usepackage[OT2,T1,T2A]{fontenc}
\usepackage[utf8x]{inputenc}
\usepackage[french,german,russian,english]{babel}
\usepackage{appendix}
\newcommand{\textcyr}[1]{
{\fontencoding{OT2}\fontfamily{wncyr}\selectfont
\hyphenchar\font=-1\relax#1}}


\DeclareMathOperator{\sech}{sech}
\DeclareMathOperator{\Span}{span}
\DeclareMathOperator{\dist}{dist}
\DeclareMathOperator{\ind}{ind}
\DeclareMathOperator{\lcm}{lcm}
\DeclareMathOperator*{\fib}{fib}

\DeclareMathOperator*{\vfib}{\widetilde{fib}}
\DeclareMathOperator{\Lyn}{Lyn}
\DeclareMathOperator{\Au}{Au}
\DeclareMathOperator{\Ag}{Ag}
\DeclareMathOperator{\disc}{disc}
\DeclareMathOperator{\Shf}{\text{\cyrins{\textsf{Ш}}}}
\DeclareMathOperator{\shf}{\text{\cyrins{\textsf{ш}}}}
\DeclareMathOperator{\Tr}{Tr}
\DeclareMathOperator{\ELi}{ELi}
\DeclareMathOperator{\supp}{supp}
\DeclareMathOperator{\Cl}{Cl}
\DeclareMathOperator{\Ls}{Ls}
\DeclareMathOperator{\ran}{ran}
\DeclareMathOperator{\sgn}{sgn}
\DeclareMathOperator{\Gal}{Gal}
\DeclareMathOperator{\End}{End}
\DeclareMathOperator*{\res}{res}
\DeclareMathOperator{\Li}{Li}
\DeclareMathOperator{\sn}{sn}
\DeclareMathOperator{\cn}{cn}
\DeclareMathOperator{\dn}{dn}
\DeclareMathOperator{\am}{am}
\DeclareMathOperator{\tn}{tn}
\DeclareMathOperator{\JZ}{Z}
\DeclareMathOperator{\Int}{Int}
\DeclareMathOperator{\arccot}{arccot}
\DeclareMathOperator{\D}{d}
\DeclareMathOperator{\I}{Im}
\DeclareMathOperator{\R}{Re}
\DeclareMathOperator*{\Reg}{Reg}

\def\RotSymbol#1#2#3{\rotatebox[origin=c]{#1}{$#2#3$}}
\def\rotcirclearrowleft{\mathpalette{\RotSymbol{-14}}\circlearrowleft}
\def\rotcurvearrowleft{\mathpalette{\RotSymbol{-100}}\curvearrowleft}
\def\turncurvearrowleft{\mathpalette{\RotSymbol{-180}}\curvearrowleft}

\def\Xint#1{\mathchoice {\XXint\displaystyle\textstyle{#1}}%
{\XXint\textstyle\scriptstyle{#1}}%
{\XXint\scriptstyle\scriptscriptstyle{#1}}%
{\XXint\scriptscriptstyle\scriptscriptstyle{#1}}%
\!\int}\def\XXint#1#2#3{{\setbox0=\hbox{$#1{#2#3}{\int}$}       \vcenter{\hbox{$#2#3$}}\kern-.5\wd0}}


\def\dashint{\Xint-}
\def\ddashint{\Xint=}
\def\clockint{\Xint\circlearrowright}
\def\counterint{\Xint\rotcirclearrowleft}
\def\pd{\text{\textit\pounds}}


%\bibpunct{[}{]}{,}{n}{}{;}

\def\qed{\hfill$ \blacksquare$}
\def\eor{\hfill$ \square$}


\def\acknowledgement{\par\addvspace{17pt}\small\rmfamily
\trivlist
\item[\hskip\labelsep
{\it{Acknowledgements}.}]}
\def\endacknowledgement{\endtrivlist\addvspace{6pt}}

\newcolumntype{L}{>{$}l<{$}}
\newcolumntype{C}{>{$}c<{$}}
\newcolumntype{R}{>{$}r<{$}}

\theoremstyle{plain}
\newtheorem{theorem}{Theorem}[section]
\newtheorem{proposition}[theorem]{Proposition}
\newtheorem{lemma}[theorem]{Lemma}
\newtheorem{corollary}[theorem]{Corollary}
\newtheorem{example}[theorem]{Example}
\newtheorem{conjecture}[theorem]{Conjecture}

\newenvironment{remark}[1][Remark]{\begin{trivlist}
\item[\hskip \labelsep {\bfseries #1}]}{\end{trivlist}}


\theoremstyle{definition}
\newtheorem{definition}[theorem]{Definition}

%\theoremstyle{remark}
%\newtheorem{notation}[theorem]{Notation}
%\newtheorem{remark}{Remark}[theorem]



\renewcommand{\thesection}{\arabic{section}}
\renewcommand*\thetable{\arabic{table}}
\numberwithin{equation}{section}

\setcounter{tocdepth}{4}


\usepackage{Baskervaldx}
\usepackage[baskervaldx]{newtxmath}
\usepackage[cal=cm,scr=rsfs,frak=euler]{mathalfa}

%\usepackage{baskervald}
%\usepackage[baskervillef]{newtxmath}

\usepackage[scr=rsfs]{mathalfa}
\DeclareMathAlphabet{\mathsf}{OT1}{\sfdefault}{m}{n}
\SetMathAlphabet{\mathsf}{bold}{OT1}{\sfdefault}{m}{n}
\DeclareSymbolFontAlphabet{\mathbb}{AMSb}
\DeclareRobustCommand{\cyrins}[1]{%
  \begingroup\fontfamily{erewhon-TLF}%
  \foreignlanguage{russian}{#1}%
  \endgroup
}

\usepackage{color}

\usepackage{eurosym}

\begin{document}

\pagenumbering{roman}
\selectlanguage{english}
%+Title
\title{Galois descents of certain multiple polylogarithms}\dedicatory{Dedicated to 19260817, a prime number}\author{Yajun Zhou}
\address{Program in Applied and Computational Mathematics (PACM), Princeton University, Princeton, NJ 08544} \email{yajunz@math.princeton.edu}\curraddr{\textrm{} \textsc{Academy of Advanced Interdisciplinary Studies (AAIS), Peking University, Beijing 100871, P. R. China}}\email{yajun.zhou.1982@pku.edu.cn}

\date{\today}\thanks{\textit{Keywords}: Galois descents, multiple polylogarithms, cyclotomic multiple zeta values\\\indent\textit{MSC 2020}: 11M06, 11M32\\\indent * This work was supported in part  by the Applied Mathematics Program within the Department of Energy
(DOE) Office of Advanced Scientific Computing Research (ASCR) as part of the Collaboratory on
Mathematics for Mesoscopic Modeling of Materials (CM4).}


%+Abstract

\begin{abstract}We prove some conjectures of K. C. Au concerning the  descents of cyclotomic levels in certain sums over multiple polylogarithms. Via answers to a question of Au in some particular cases, we also confirm Broadhurst's conjecture on honorary multiple zeta values at even weights. \end{abstract}
%-Abstract

\maketitle
%+Contents
\pagenumbering{arabic}

\section{Introduction\label{sec:intro}}
Multiple polylogarithms (MPLs)  are defined by (analytic continuations of) \begin{align}
\Li_{a_1,\dots,a_n}(z_1,\dots,z_n)\colonequals {\sum\limits_{\ell_{1}>\dots>\ell_{n}>0}\prod\limits_{j=1}^n\frac{\smash[t]{z_{j}^{\ell_{j}}}}{\ell_j^{a_j}}}.\label{eq:Mpl_defn}
\end{align} Here, the right-hand side of \eqref{eq:Mpl_defn}  converges absolutely  for  $ a_1,\dots,a_n\in\mathbb Z_{>0}, \prod_{j=1}^m|z_j|<1,m\in\mathbb Z\cap[1,n]$. The series in question also converges (though not necessarily in the mode of absolute convergence) when  $ (a_1,z_1)\neq(1,1)$ and $ |z_1|=\dots=|z_n|=1$. A convergent MPL $ \Li_{a_1,\dots,a_n}(z_1,\dots,z_n)$ is said to have weight $ a_1+\dots+a_n$ and depth $n$.

Using convergent MPLs, one may introduce the $ \mathbb Q$-vector space spanned by cyclotomic multiple zeta values (CMZVs) of weight $ w\in\mathbb Z_{>0}$ and level $N\in\mathbb Z_{>0}$:\begin{align}
 \mathfrak Z_{w}(N)\colonequals\Span_{\mathbb Q}\left\{\Li_{a_1,\dots,a_n}(z_1,\dots,z_n)\colonequals \smash[b]{\sum\limits_{\ell_{1}>\dots>\ell_{n}>0}\prod\limits_{j=1}^n\frac{\smash[t]{z_{j}^{\ell_{j}}}}{\ell_j^{a_j}}}\middle|\begin{smallmatrix}a_1,\dots,a_n\in\mathbb Z_{>0}\\z_{1}^{N}=\dots=z_n^N=1\\\sum _{j=1}^{n}a_{j}=w\\(a_1,z_1)\neq(1,1)\end{smallmatrix} \right\}.\label{eq:Zk(N)}
\end{align}Conventionally, one sets $ \mathfrak Z_{0}(N) \colonequals \mathbb Q$. CMZVs of levels 1 and 2 are known as multiple zeta values (MZVs) and alternating multiple zeta values (AMZVs), respectively. Customarily, (A)MZVs are  denoted by \begin{align}
\zeta_{A_{1},\dots,A_n}\equiv\Li_{|A_1|,\dots,|A_n|}\bigg(\frac{A_1}{|A_1|},\dots,\frac{A_n}{|A_n|}\bigg)\in \mathfrak Z_{w}(2),
\end{align}where $ A_1,\dots,A_n\in\mathbb Z\smallsetminus\{0\}$ and $ A_1\neq1$. With symbolic computations involving provable algebraic relations\footnote{In particular, Au's software \cite{Au2022a} gives automated reductions  in $ \mathfrak Z_{w}(N)$ for $ w\in\{1,2,3,4,5,6\}$, $N=4$ and  $ w\in\{1,2,3,4,5\}$, $ N=6$.} among CMZVs of mod\-er\-ate-sized weights \cite{Au2022a}, Kam Cheong Au discovered some remarkable patterns concerning \begin{align}\Li_{a_1,\dots,a_n}(z,\mathbf 1_{n-1})
\equiv\Li_{a_1,\dots,a_n}(z,\underset{n-1}{\underbrace{1,\dots,1}})\colonequals \sum_{\ell_{1}>\dots>\ell_{n}>0}\frac{z^{\ell_{1}}}{\prod_{j=1}^n\ell_j^{a_j}}\quad(\text{where } |z|\leq1)
\end{align}and \begin{align}
g_{w,n}(z)\colonequals {}&(-1)^n\sum_{\substack{a_1,\dots,a_n\in\mathbb Z_{>0}\\a_1+\dots+a_n=w}} \Li_{a_1,\dots,a_n}(z,\mathbf 1_{n-1}),\label{eq:g_w,n}
\end{align}
as recapitulated in the conjecture below.

\begin{conjecture}[K. C. Au {\cite[Conjecture 1.5]{Au2022a}}]\label{conj:Au1.5}\begin{enumerate}[leftmargin=*,  label=\emph{(\alph*)},ref=(\alph*),
widest=d, align=left] \item
For each fixed pair of weight $w\in\mathbb Z_{>0}$ and depth $ n\in\mathbb Z\cap[1,w]$, one has  {\allowdisplaybreaks
\begin{align}
\R g_{w,n}\bigg( \frac{1+i}{2} \bigg)\in{}&\mathfrak Z_w(2),\label{eq:Au1.5a1}\\\R g_{w,n}( i )\in{}&\mathfrak Z_w(2),\label{eq:Au1.5a2}
\end{align}
}where each individual summand is in $ \mathfrak Z_w(4)$.\footnote{In an unlabeled formula right above  \cite[Conjecture 1.5]{Au2022a}, Au stated a borderline case of \eqref{eq:Au1.5a1}, where $n=1$.}
\item Set $ \varrho\colonequals  e^{\pi i/3}$. For each fixed pair of weight $w\in\mathbb Z_{>0}$ and depth $ n\in\mathbb Z\cap[1,w]$, one has \begin{align}
\R g_{w,n}(\varrho)\in{}&\mathfrak Z_w(1),\label{eq:Au1.5b}
\end{align}where each individual summand is in $ \mathfrak Z_w(6)$.\footnote{\label{fn:desc3}According to either \cite[Example 3.9]{Au2022a} or \cite[Proposition 3.3]{SunZhou2026MCV}, there is a sharper result $ \mu_{a_1,\dots, a_n}\in\mathfrak Z_{a_1+\dots+a_n}(3)$.  Also note that one may verify \eqref{eq:Au1.5b} for $ w\in\mathbb Z\cap[1,8]$ and $ n\in\mathbb Z\cap[1,w]$, using the algorithms given in \cite[\S\S3.1--3.2]{SunZhou2026MCV}. }
 \end{enumerate}
\end{conjecture}

Au referred to these  relations as ``Galois descent for multiple polylogarithms---a subtle and
largely unexplored phenomenon''.
 Here, 
 {the   wording  ``largely unexplored'' was particularly  judicious. Galois descents of CMZVs did not belong to an uncharted territory before Au's work~\cite{Au2022a}: Broadhurst's numerical experiments in 2014 suggested the presence of honorary multiple zeta values  (or just ``hononary MZVs'' for short) at even weights  \cite[Conjecture 5]{Broadhurst2014MDV}\begin{align}
\R \Li_{w-1,1}(\varrho,1)\in\mathfrak Z_w(1),\quad \text{for }w\in 2\mathbb Z_{>0},\label{eq:BroadhurstConj5}
\end{align} while  Xu's empirical discovery  \cite[Conjecture 7.2]{Xu2019nonlinEuler}  \begin{align}\begin{split}&\zeta_{-2m,2}+2m\zeta_{-2m-1,1}\equiv\Li_{2m,2}(-1,1)+2m\Li_{2m+1,1}(-1,1)\\\in{}& \mathfrak Z_{2m+2}(1)\cap\Span_{\mathbb Q}\{\zeta_{a,2m+2-a}\equiv\Li_{a,2m+2-a}(1,1)|a\in\mathbb Z\cap[1,2m+1]\}\label{eq:XuConj}
\end{split}\end{align}  in 2019 has been  proved very recently by Zheng--Yang  \cite{ZhengYang2025}  (in part) and by Charlton  \cite{Charlton2026} (in full). Moreover, by working with the categories of mixed Tate motives, Glanois \cite{Glanois2016} established criteria for Galois descents in motivic versions for CMZVs of levels $N\in\{ 2,3,4,6,8\}$, which paved the way for the proofs of Zheng--Yang~\cite{ZhengYang2025} and  Charlton~\cite{Charlton2026}.}

In this note, we not only answer Au's questions \eqref{eq:Au1.5a1}--\eqref{eq:Au1.5b} in the positive, but also confirm Broadhurst's \eqref{eq:BroadhurstConj5}.

We open with a brief  overview of our  toolkit in  \S\ref{sec:MPL_GPL}, before establishing ``depth-blind'' versions of \eqref{eq:Au1.5a1} and \eqref{eq:Au1.5b}, which replace their left-hand sides by (alternating) sums over $ n\in\mathbb Z\cap[1,w]$.  In \S\ref{sec:desc3}, we prove the full versions of  both \eqref{eq:Au1.5b} and \eqref{eq:BroadhurstConj5},  upon a conversion of $ g_{w,n}(z)$ from its defining multi-index summation [see \eqref{eq:g_w,n}] to an integral representation.
Afterwards, these techniques are extended to the complete verifications of \eqref{eq:Au1.5a1}  and \eqref{eq:Au1.5a2} in \S\ref{sec:desc4}, along with another conjecture of Au \cite[Conjecture 1.6]{Au2022a}. \section{ Depth-blind Galois descents of multiple polylogarithms \label{sec:MPL_GPL}}
\subsection{Generalized polylogarithms and their properties\label{subsec:GPLprop}}
The generalized polylogarithms (GPLs) are   defined through the recursion \cite[(2.1)]{Frellesvig2016} \begin{align}\notag\\[-12pt]
G(\alpha_{1},\dots,\alpha_n;z)\colonequals{\displaystyle\int_0^z\frac{G(\alpha_2,\dots,\alpha_n;x)\D x}{x-\alpha_1}},\quad\text{for }\displaystyle\smash[t]{\sum_{k=1}^n|}\alpha_k|\neq0\label{eq:GPL_rec}
\end{align}and the initial condition \cite[(2.2)]{Frellesvig2016}\begin{align}
{ G(\boldsymbol0_{n};z)\equiv G(\underset{n }{\underbrace{0,\dots,0 }};z)}\colonequals{\dfrac{\log^nz}{n!}},\quad G(-\!\!-;z)\colonequals1.\label{eq:GPL0}
\end{align}To better appreciate the combinatorial structure of MPLs, we convert them to GPLs  (see \cite[(2.5)]{Frellesvig2016} or \cite[(1.3)]{Panzer2015}): \begin{align}
\begin{split}\Li_{a_1,\dots,a_n}(z_1,\dots,z_n)={}&(-1)^{n}G\left(\smash[b]{\underset{a_1-1 }{\underbrace{0,\dots,0 }}},\frac{1}{z_{1}},\smash[b]{\underset{a_2-1 }{\underbrace{0,\dots,0 }}},\frac{1}{z_{1}z_2},\dots,\smash[b]{\underset{a_n-1 }{\underbrace{0,\dots,0 }}},\frac{1}{\prod_{j=1}^nz_j};1\right)\\[-5pt]\\\equiv{}&(-1)^nG\left(  \boldsymbol0_{a_{1}-1},\frac{1}{z_{1}},\boldsymbol0_{a_{2}-1},\frac{1}{z_{1}z_2},\dots,\boldsymbol0_{a_{n}-1},\frac{1}{\prod_{j=1}^nz_j};1\right).
\end{split}\label{eq:MPLtoGPL}
\end{align}Note that when  $ \alpha_n\neq0$, the GPL recursion \eqref{eq:GPL_rec} effectively defines $G(\alpha_{1},\dots,\alpha_n;z)$ as a function of the homogeneous coordinates $ [\alpha_{1}:\cdots:\alpha_n:z]$, namely  \cite[(2.3)]{Frellesvig2016}
\begin{align}
G(\alpha_{1},\dots,\alpha_n;z)=G(\sigma\alpha_{1},\dots,\sigma\alpha_n;\sigma z)\label{eq:GPLscaling}
\end{align}for $\sigma\neq0 $.
In particular, we have\begin{align}
\begin{split}\begin{split}\Li_{a_1,\dots,a_n}(z,\mathbf 1_{n-1})={}&(-1)^{n}G({\underset{a_1-1 }{\underbrace{0,\dots,0 }}},1,{\underset{a_2-1 }{\underbrace{0,\dots,0 }}},1,\dots,\smash[b]{\underset{a_n-1 }{\underbrace{0,\dots,0 }}},1;z)\\\equiv{}&(-1)^{n}G(\boldsymbol0_{a_1-1},1,\boldsymbol0_{a_2-1},1,\dots,\boldsymbol0_{a_n-1},1;z),
\end{split}\label{eq:LtoGPL}
\end{split}
\end{align}according to the MPL-GPL conversion \eqref{eq:MPLtoGPL}  and the scaling property \eqref{eq:GPLscaling}.

 The products of GPLs satisfy \cite[(2.4)]{Frellesvig2016} \begin{align}\label{eq:GPL_shuffle}G(\alpha_{1},\dots,\alpha_{j};z)G(\beta_{1},\dots,\beta_{k};z)=\sum_{\bm\gamma\in\bm \alpha\Shf\bm \beta }G(\gamma_{1},\dots,\gamma_{j+k}; z)\end{align}where the set $ \bm \alpha\Shf\bm \beta$ runs over all the shuffles of the vector components in  $ \bm \alpha$ and $\bm \beta$   that maintain the internal orders  within both vectors. 

\subsection{Equations of motion for  ``depth-blind'' Galois descents\label{subsec:depth-blind}}
In view of \eqref{eq:LtoGPL}, we can rewrite \eqref{eq:g_w,n} as\begin{align}
g_{w,n}(z)\colonequals {}&\sum_{\substack{\alpha_1,\dots,\alpha_{w-1}\in\{0,1\}\\\alpha_1+\dots+\alpha_{w-1}+1=n}}G(\alpha_1,\dots,\alpha_{w-1},1;z) .\tag{\ref{eq:g_w,n}$'$}
\end{align}
Now, we settle ``depth-blind'' versions of \eqref{eq:Au1.5a1} and \eqref{eq:Au1.5b}
for all the positive weights, as described in the two propositions below.\begin{proposition}

\label{prop:G_avg_desc}Define\begin{align}
\mathsf G_w(z)\colonequals {}&\sum_{n=1}^w g_{w,n}(z)
\end{align}for all $ w\in\mathbb Z_{>0}$. We have \begin{align}
\mathsf G_w(z)+\mathsf G_w(1-z)\in\Span_\mathbb Q\left\{ Z_\ell [\log z+\log(1-z)]^{w-\ell}\middle|\begin{smallmatrix}\ell,w-\ell\in\mathbb Z_{\geq0}\\Z_\ell\in\mathfrak Z_\ell(1)\end{smallmatrix}\right\}.\label{eq:Gw_refl}
\end{align}In particular, this implies that \begin{align}
\R \mathsf G_w\bigg(\frac{1+i}{2}\bigg)\in{}&\mathfrak Z_w(2),\\\R \mathsf G_{w}(\varrho)\in{}&\mathfrak Z_w(1).
\end{align}
\end{proposition}
\begin{proof}Thanks to the GPL recursion \eqref{eq:GPL_rec}, we obtain an equation of motion\footnote{Since empty sums are zero, we   have $ g_{w,n}(z)=0$ for $ n\notin\mathbb Z\cap[1,w]$. }\begin{align}
\frac{\partial g_{w,n}(z)}{\partial z}=\frac{ g_{w-1,n}(z)}{z}+\frac{g_{w-1,n-1}(z)}{z-1}.\label{eq:gwn_evolv}
\end{align}Summing over all the depths $ n\in\mathbb Z\cap[1,w]$, we get \begin{align}
\frac{\partial\mathsf G_w(z)}{\partial z}={}&\left( \frac{1}{z} +\frac{1}{z-1}\right)\mathsf G_{w-1}(z).
\end{align}Trading $z$ for $1-z$, one has\begin{align}
\frac{\partial\mathsf G_w(1-z)}{\partial z}={}&\left( \frac{1}{z} +\frac{1}{z-1}\right)\mathsf G_{w-1}(1-z).
\end{align}To deduce \eqref{eq:Gw_refl} from the recursion \begin{align}
\frac{\partial[\mathsf G_w(z)+\mathsf G_w(1-z)]}{\partial z}={}&\left( \frac{1}{z} +\frac{1}{z-1}\right)[\mathsf G_{w-1}(z)+G_{w-1}(1-z)]
\end{align}with initial condition\begin{align}
\mathsf G_1(z)+\mathsf G_1(1-z)=\log z+\log(1-z),
\end{align}one only needs to check the regularized limits \cite[\S2.3]{Panzer2015}\begin{align}
\Reg_{z\to0}[\mathsf G_w(z)+\mathsf G_w(1-z)]=\Reg_{z\to0}\mathsf G_w(1-z)\in \mathfrak Z_w(1),\label{eq:RegG(1-z)}
\end{align}which is a routine exercise.
[By working a little harder, one can actually represent $ \Reg_{z\to0}\mathsf G_w(1-z)$ as a polynomial (in rational coefficients) of $ \pi^2$ and all the odd zeta values $ \zeta_{2n+1}\equiv\Li_{2n+1}(1),n\in\mathbb Z_{>0}$. See  \S\ref{subsec:shf_gwn} for details.]

The rest  follow from the observations that $ \log\frac{1+i}{2}+\log\frac{1-i}{2}=\Li_1(-1)\in\mathfrak Z_1(2)$ and $ \log \varrho +\log(1-\varrho)=0$, as well as the GPL shuffles in \eqref{eq:GPL_shuffle}.\end{proof}\begin{proposition}\label{prop:G_alt_avg_desc}Define\begin{align}
\widetilde{\mathsf G}_w(z)\colonequals {}&\sum_{n=1}^w (-1)^{n}g_{w,n}(z)
\end{align}for all $ w\in\mathbb Z_{>0}$. For $ |z|\le1$ and $ z\neq1$, we have
\begin{align}
\widetilde{\mathsf G}_w(z)= g_{w,1}\bigg(\frac{z}{z-1}\bigg)\equiv-\Li_w\bigg(\frac{z}{z-1}\bigg),
\label{eq:gw1Landen}\end{align}which entails \begin{align}
\R \widetilde{\mathsf G}_1\bigg(\frac{1+i}{2}\bigg)=\frac{\log2}{2}\in\mathfrak Z_1(2),\quad \R \widetilde {\mathsf G}_1(\varrho)=0\in\mathfrak Z_1(1)=\{0\}
\end{align}and\begin{align}
\R \widetilde{\mathsf G}_w\bigg(\frac{1+i}{2}\bigg)\in \mathbb Q\zeta _{w}\subseteq\mathfrak Z_w(1),\\\R \widetilde{\mathsf G}_w(\varrho)\in \mathbb Q\zeta _{w}\subseteq\mathfrak Z_w(1),
\end{align}for all $ w\in\mathbb Z_{>1}$.
\end{proposition}\begin{proof}The function\begin{align}
\widetilde{\mathsf \Delta}_w(z)\colonequals\widetilde{\mathsf G}_w(z)-g_{w,1}\bigg(\frac{z}{z-1}\bigg) \end{align}is governed by the recursion [cf.\ \eqref{eq:gwn_evolv}] \begin{align}
\frac{\partial \widetilde{\mathsf \Delta}_w(z)}{\partial z}=\left( \frac{1}{z} -\frac{1}{z-1}\right)\widetilde{\mathsf \Delta}_{w-1}(z)
\end{align}and the boundary condition $ \widetilde{\mathsf\Delta}_w(0)=0$ for all  $ w\in\mathbb Z_{>1}$. In other words, one has\begin{align}
\widetilde{\mathsf \Delta}_w(z)=\int_{0}^z\left( \frac{1}{x} -\frac{1}{x-1}\right)\widetilde{\mathsf \Delta}_{w-1}(x)\D x
\end{align} for all  $ w\in\mathbb Z_{>1}$. Direct computation reveals that  $ \widetilde{\mathsf\Delta}_1(z)=0$, so we must have  $ \widetilde{\mathsf \Delta}_w(z)=0$ for all $ z\in\mathbb C$ and   $ w\in\mathbb Z_{>1}$.

To wrap up this proposition, simply note that  \begin{align}
\R\Li_w(-i)={}&\frac{1-2^{w-1}}{2^{2w-1}}\zeta _{w}\intertext{and}\R\Li_w(\varrho)={}&\left( 1-\frac{1}{2^{w-1}} \right)\left( 1-\frac{1}{3^{w-1}} \right)\frac{\zeta _{w}}{2}
\end{align}hold for all $ w\in\mathbb Z_{>1}$.
\end{proof}

\section{ Galois descents from hexagon vertices\label{sec:desc3}}
\subsection{Shuffle structure of Au's sums\label{subsec:shf_gwn}}The  ``depth-blind'' appetizers in \S\ref{subsec:depth-blind} primarily drew on the GPL recursion  \eqref{eq:GPL_rec},  without making heavy use of the GPL shuffles in  \eqref{eq:GPL_shuffle}.  In this subsection, we will produce identities of a different flavor, to investigate the ``in-depth'' structure of Au's sum rules.

We start from an application of the  GPL shuffles in  \eqref{eq:GPL_shuffle}:\begin{align}
G(1;z)G(\boldsymbol0_{w-1},1;z)=G(\boldsymbol0_{w-1},1,1;z)+g_{w,2}(z),\label{eq:gw2_prod}
\end{align}  or in terms of word shuffles,\begin{align}
1\shf \boldsymbol0_{w-1}1= \boldsymbol0_{w-1}11+\sum_{\substack{\alpha_1,\dots,\alpha_{w-1}\in\{0,1\}\\\alpha_1+\dots+\alpha_{w-1}+1=2}}\alpha_1\cdots\alpha_{w-1}1.
\tag{\ref{eq:gw2_prod}$'$}\label{eq:gw2_prod'}\end{align}   In the next lemma, we will generalize \eqref{eq:gw2_prod} and \eqref{eq:gw2_prod'} to all the available weights and depths. \begin{lemma}
For $ w\in\mathbb Z_{>0}$ and $ n\in\mathbb Z\cap[1,w]$, we have \begin{align}
g_{w,n}(z)=\sum_{m=1}^n(-1)^{m-1}G(\boldsymbol1_{n-m};z)G(\boldsymbol0_{w-n},\boldsymbol1_{m};z),\label{eq:gwn_prod_expn}
\end{align}with the understanding that $G(\boldsymbol0_{0},\boldsymbol1_{m};z)=G(\boldsymbol1_{m};z)$ and $ G(\boldsymbol1_{0};z)=1$. \end{lemma}\begin{proof}For $ n=1$,  the proposed formula follows from the definition of $ g_{w,1}(z)$.  For $n=w$, we have \begin{align}
g_{w,w}(z)=G(\boldsymbol1_{w};z)=\frac{\log^w(1-z)}{w!}
\end{align}by  iterative invocations of either the GPL recursion  \eqref{eq:GPL_rec} or the GPL shuffles   \eqref{eq:GPL_shuffle} for $ 1\shf\textbf{1}_n=(n+1)\boldsymbol1_{n+1}$, whereas the right-hand side of \eqref{eq:gwn_prod_expn} becomes \begin{align}
\sum_{m=1}^w(-1)^{m-1}G(\boldsymbol1_{ w-m};z)G(\boldsymbol1_{m};z)=\sum_{m=1}^w(-1)^{m-1}\frac{\log^w(1-z)}{(w-m)!m!}=\frac{\log^w(1-z)}{w!},
\end{align}according to the binomial theorem.

 In the light of \begin{align}
\frac{\partial G(\boldsymbol1_{n-m};z)}{\partial z}={}&\begin{cases}0 & \text{if }n-m=0, \\
\frac{ G(\boldsymbol1_{n-m-1};z)}{z-1} & \text{if }n-m>0, \\
\end{cases}\intertext{and}\frac{\partial G(\boldsymbol0_{w-n},\boldsymbol1_{m};z)}{\partial z}={}&\begin{cases}\frac{G(\boldsymbol1_{m-1};z)}{z-1} & \text{if }w-n=0, \\
\frac{ G(\boldsymbol0_{w-n-1},\boldsymbol1_{m};z)}{z} & \text{if }w-n>0, \\
\end{cases}
\end{align}we see that\begin{align}
\delta_{w,n}(z)\colonequals g_{w,n}(z)-\sum_{m=1}^n(-1)^{m-1}G(\boldsymbol1_{n-m};z)G(\boldsymbol0_{w-n},\boldsymbol1_{m};z)
\end{align}evolves in the same fashion as $  g_{w,n}(z)$:\begin{align}
\frac{\partial \delta_{w,n}(z)}{\partial z}=\frac{ \delta_{w-1,n}(z)}{z}+\frac{\delta_{w-1,n-1}(z)}{z-1},\quad\text{where }w\in\mathbb Z_{>0},n\in\mathbb Z\cap[1,w],
\end{align}so long as one adopts the convention that $ \delta_{w-1,w}(z)=0$. Since $ \delta_{w,n}(0)=0$, we may evaluate\begin{align}
 \delta_{w,n}(z)=\int_{0}^z\left[ \frac{ \delta_{w-1,n}(x)}{x}+\frac{\delta_{w-1,n-1}(x)}{x-1} \right]\D x\label{eq:delta_int_evolv}
\end{align} iteratively,  as follows:\begin{itemize}
\item
Using $ \delta_{2,1}(x)=\delta_{2,2}(x)=0$ proved in the last paragraph, deduce $ \delta_{3,2}(z)=0$.\item Fixing the depth $n=2$, prove $ \delta_{w,2}(z)=0$  by induction on the weight  $w$, through repeated invocations of \eqref{eq:delta_int_evolv}. [Here, one does not need to know the shuffle identity \eqref{eq:gw2_prod} beforehand.]\item Using $ \delta_{3,3}(x)=0$ proved in the last paragraph, deduce $ \delta_{4,3}(z)=0$, and then promote to all the weights $ w>4$ while fixing the depth  $n=3$.\item Proceed as above, advancing the depth $n$ by one step at a time, before validating $ \delta_{w,n}(z)=0$ for all $ w\geq n$.
\end{itemize} This completes the proof of  \eqref{eq:gwn_prod_expn}.
\end{proof}\begin{remark}In terms of word shuffles, the last lemma converts a multi-index summation over certain binary words  into a single sum over $ m\in\mathbb Z\cap[1,n]$:\begin{align}
\sum_{\substack{\alpha_1,\dots,\alpha_{w-1}\in\{0,1\}\\\alpha_1+\dots+\alpha_{w-1}+1=n}}\alpha_1\cdots\alpha_{w-1}1= \sum_{m=1}^n(-1)^{m-1}\boldsymbol1_{n-m}\shf\boldsymbol0_{w-n}\boldsymbol1_{m},\tag{\ref{eq:gwn_prod_expn}$'$}\label{eq:gwn_prod_expn'}
\end{align}so long as  the empty words $ \boldsymbol0_0$ and $ \boldsymbol1_0$  are suitably interpreted.
\eor\end{remark}\begin{remark}During the proof of Proposition \ref{prop:G_avg_desc}, we encountered a quantity\begin{align}
\Reg_{z\to0}\mathsf G_w(1-z)=\Reg_{z\to0} \sum _{n=1}^wg_{w,n}(1-z),
\end{align} where the operation    ``$ \Reg_{z\to0}$'' was meant to suppress all the positive integer powers of $ \log z$. With \eqref{eq:gwn_prod_expn'}, we have\begin{align}
\begin{split}\Reg_{z\to0} g_{w,n}(1-z)={}&\Reg_{z\to0}\sum_{m=1}^n(-1)^{m-1}G(\boldsymbol1_{n-m};1-z)G(\boldsymbol0_{w-n},\boldsymbol1_{m};1-z)\\={}&(-1)^{n-1}G(\boldsymbol0_{w-n},\boldsymbol1_{n};1)
\end{split}
\end{align}for $ n\in\mathbb Z\cap[1,w-1]$ and $ \Reg_{z\to0}g_{w,w}(1-z)=\frac{1}{w!}\Reg_{z\to0}\log^w z=0$, so \begin{align}
\begin{split}\Reg_{z\to0}\mathsf G_w(1-z)={}&\sum_{n=1}^{w-1}(-1)^{n-1}G( \boldsymbol0_{w-n},\boldsymbol1_{n};1)=-\sum_{n=1}^{w-1}\Li_{w-n+1,\boldsymbol1_{n-1}}(\boldsymbol1_{n})\\\underset{\text{\cite[(5)]{BorweinStraub2015Snp}}}{\xlongequal{\text{\cite[(1.3)]{Koelbig1986}}}}{}&-\sum_{n=1}^{w-1}\frac{(-1)^{w-1}}{(w-n-1)!n!}\int_0^1\frac{\log^{w-n-1}t\log^{n}(1-t)\D t}{t}\\={}&\frac{(-1)^w}{(w-1)!}\int_0^1\frac{[\log t+\log(1-t)]^{w-1}-\log^{w-1}t}{t}\D t\\\in{}&\mathfrak Z_{w}(1)\cap\mathbb Q[\pi^2,\zeta_{3},\zeta_{5},\dots]
\end{split}\label{eq:RegGw(1-z)}
\end{align} gives a quantitative refinement of the statement in \eqref{eq:RegG(1-z)}, where  the ``$ \in$'' step follows from a standard argument \cite[\S9.1]{Koelbig1986}.
\eor\end{remark}

To facilitate further discussions, we rewrite the shuffle identities \eqref{eq:gwn_prod_expn} and \eqref{eq:gwn_prod_expn'}
with the  integral representations of Nielsen polylogarithms, as studied by K\"olbig \cite[(1.3)]{Koelbig1986} and Borwein--Straub \cite[(5)]{BorweinStraub2015Snp}. \begin{proposition}
For $ |z|\leq1$, $ z\neq1$, $ w\in\mathbb Z_{>0}$ and $ n\in\mathbb Z\cap[1,w-1]$, we have\begin{align}
g_{w,n}(z)=\frac{(-1)^w}{(w-n-1)!n!}\int_0^1 \frac{[\log^{}(1-zt)-\log(1-z)]^{n}-[-\log(1-z)]^n}{t}\log^{w-n-1}t\D t.\label{eq:gwn_int}
\end{align} \end{proposition}\begin{proof}We compute \begin{align}
\begin{split}g_{w,n}(z)={}&-\sum_{m=1}^n\frac{\log^{n-m}(1-z)}{(n-m)!}\Li_{w-n+1,\boldsymbol1_{m-1}}(z,\boldsymbol1_{m-1})\\\underset{\text{\cite[(5)]{BorweinStraub2015Snp}}}{\xlongequal{\text{\cite[(1.3)]{Koelbig1986}}}}{}&\frac{1}{(w-n-1)!}\sum_{m=1}^n\frac{(-1)^{w-n+m}\log^{n-m}(1-z)}{(n-m)!m!}\int_0^1\frac{\log^{w-n-1}t\log^{m}(1-zt)\D t}{t}.
\end{split}
\end{align}  The binomial theorem then leads us to  \eqref{eq:gwn_int}.
\end{proof}
\begin{remark}The integral representations for Nielsen polylogarithms do not work when $n=w$. Fortunately,  direct computations reveal these exceptional cases as simply $ g_{w,w}(z)=\frac{\log^w(1-z)}{w!}$ for all $ w\in\mathbb Z_{>0}$.
\eor\end{remark}

Consider the hypergeometric  series\begin{align}
_2F_1\left( \left.\begin{array}{@{}c@{}}
a,b \\
c \\
\end{array}\right|z \right)=1+\sum_{k=1}^\infty\frac{(a)_k(b)_k}{(c)_k}\frac{z^k}{k!},\quad|z|\leq 1,z\neq1, \label{eq:2F1_defn}
\end{align}
where $ (A)_{k}=\prod_{m=0}^{k-1}(A+m)$ is the rising factorial. In the next proposition, we turn the right-hand side of    \eqref{eq:gwn_int} into partial derivatives of  $ _2F_1$   series with respect to their  hypergeometric parameters.\begin{proposition}
For $ |z|\leq1$, $ \left\vert \frac{z}{z-1} \right\vert\leq1$, $ z\neq1$, and $ m+1,n\in\mathbb Z_{>0}$, we have\begin{subequations}\begin{align}
g_{m+n,n}(z)={}&\frac1{m!n!}\left.\!\frac{\partial^{m+n}}{\partial u^{m}\partial v^n}\left[ (1-z)^{v}{_2}F_1\left(\begin{array}{@{}c@{}}
-u,v \\
1-u \\
\end{array}\middle| z\right) -1\right]\right|_{u=v=0}\label{eq:gwn_dd_2F1}\\={}&\frac1{m!n!}\left.\!\frac{\partial^{m+n}}{\partial u^{m}\partial v^n}\left[ {_2}F_1\left(\begin{array}{@{}c@{}}
1,v \\
1-u \\
\end{array}\middle| \frac{z}{z-1}\right) -1\right]\right|_{u=v=0}.\label{eq:gwn_dd_2F1Pfaff}
\end{align}\end{subequations}\end{proposition}\begin{proof}With Euler's integral representation for the hypergeometric series  $ _2F_1$ \cite[Theorem 2.2.1]{AAR}, we have\begin{align}
\begin{split}\sum _{m=0}^\infty \sum _{n=1}^\infty g_{m+n,n}(z)u^mv^n={}&\sum _{n=1}^\infty \frac{\log^n(1-z)}{n!}v^n+u\int_0^1 \frac{(1-z)^{v}-\left(\frac{1-z}{1-zt}\right)^v}{t^{u+1}}\D t\\={}&[(1-z)^{v}-1]+(1-z)^{v}\left[{_2}F_1\left(\begin{array}{@{}c@{}}
-u,v \\
1-u \\
\end{array}\middle| z\right) -1 \right]\\={}&(1-z)^{v}{_2}F_1\left(\begin{array}{@{}c@{}}
-u,v \\
1-u \\
\end{array}\middle| z\right)-1
\end{split}
\end{align}for sufficiently small $|u| $ and $|v|$, so  \eqref{eq:gwn_dd_2F1} follows immediately.

Pfaff's transformation \cite[(2.2.6)]{AAR}  brings us\begin{align}
{_2}F_1\left(\begin{array}{@{}c@{}}
v,-u \\
1-u \\
\end{array}\middle| z\right)=(1-z)^{-v}{_2}F_1\left(\begin{array}{@{}c@{}}
v,1 \\
1-u \\
\end{array}\middle| \frac{z}{z-1}\right),
\end{align} which explains \eqref{eq:gwn_dd_2F1Pfaff}.
\end{proof}


 
\subsection{Au's Galois descents from ``level 6'' to level 1\label{subsec:desc3}}
Now we embark on the proof of  Au's  sum rule in \eqref{eq:Au1.5b}.\footnote{As indicated in Footnote \ref{fn:desc3}, for each given weight $ w$ and level $n$, this sum rule represents a Galois descent from  $ \mathfrak Z_w(3)$ [in the guise of $ \mathfrak Z_w(6)$, hence the presence of quotation marks in the title of this subsection] to   $ \mathfrak Z_w(1)$. }  Towards this end, we unpack\begin{align}\R
g_{m+n,n}(\varrho)=\left.\!\frac1{m!n!}\R\frac{\partial^{m+n}}{\partial u^{m}\partial v^n}\left[ {_2}F_1\left(\begin{array}{@{}c@{}}
1,v \\
1-u \\
\end{array}\middle| \frac{1}{\varrho}\right) -1\right]\right|_{u=v=0}
\end{align}into Euler's integral representation \cite[Theorem 2.2.1]{AAR}
as follows:\begin{align}
\begin{split}\R
g_{m+n,n}(\varrho)={}&\frac1{m!n!}\left.\!\R\frac{\partial^{m+n}}{\partial u^{m}\partial v^n}\left[ \frac{\Gamma (1-u) }{\Gamma (v) \Gamma (1-u-v)} \int_0^1 \frac{t^{v-1} (1-t)^{-u-v}}{1-\frac{t }{\varrho}} \D t-1\right]\right|_{u=v=0}\\={}&\frac1{m!n!}\left.\!\R\frac{\partial^{m+n}}{\partial u^{m}\partial v^n}\left[ \frac{1 }{\mathrm B(1-u-v,v)} \int_0^1 \frac{t^{v-1} (1-t)^{-u-v}}{1-\frac{t }{\varrho}} \D t-1\right]\right|_{u=v=0},
\end{split}
\end{align}where $ \mathrm B(\cdot,\cdot)$ 
and $ \Gamma(\cdot)$ are Euler's beta and gamma functions. Now that\begin{align}
\left.\!\frac{\partial^{m+n}}{\partial u^{m}\partial v^n}\left[ \frac{1 }{\mathrm B(1-u-v,v)} -v\right]\right|_{u=v=0}\in \mathfrak Z_{m+n-1}(1)
\end{align}is expressible through the Riemann zeta values $ \zeta_{w}\equiv\Li_w(1)$ for $ w\in\mathbb Z_{>1}$ (cf.\ \cite[(2.17)]{Xu2017}), it will suffice to show that \begin{align}
\left.\!\R\frac{\partial^{m+n}}{\partial u^{m}\partial v^n}\int_0^1 \frac{t^{v-1} (1-t)^{-u-v}}{1-\frac{t }{\varrho}} \D t\right|_{u=v=0}
\end{align}are representable through MZVs for $ m,n\in\mathbb Z_{>0}$. Furthermore, since\begin{align}
\R\frac{1}{\left(1-\frac{t }{\varrho} \right)t}=\frac{1}{t}-\frac{1}{2}\left( \frac{1}{t-\varrho} +\frac{1}{t-\frac{1}{\varrho}}\right),
\end{align}and\begin{align}
\int_0^1\frac{\log^{a} t \log^{b}(1-t)}{t}\D t\in\mathfrak Z_{a+b+1}(1) 
\end{align}is a consequence of GPL recursion \eqref{eq:GPL_rec}  and GPL shuffles \eqref{eq:GPL_shuffle} for all $a,b\in\mathbb Z_{>0}$, we may focus on the rational function\begin{align}
\frac{1}{t-\varrho} +\frac{1}{t-\frac{1}{\varrho}}=\frac{2t-1}{1-t(1-t)}
\end{align}whose poles are two vertices of a hexagon inscribed in the unit circle.\begin{theorem}\label{thm:1.5b}
We have \begin{align}
\int_0^1\left( \frac{1}{t-\varrho} +\frac{1}{t-\frac{1}{\varrho}} \right)\log^{a} t \log^{b}(1-t)\D t\in\mathfrak Z_{a+b+1}(1)\label{eq:desc3int}
\end{align}for all $a,b\in\mathbb Z_{>0}$. As a result, Au's relation $\R g_{w,n}(\varrho)\in\mathfrak Z_w(1)$ holds for all $ w\in\mathbb Z_{>0}$ and $ n\in\mathbb Z\cap[1,w]$.\end{theorem}
\begin{proof}We only need to work out \begin{align}
   \mathscr I_{a,b}\colonequals\int_0^1\left( \frac{1}{t-\varrho} +\frac{1}{t-\frac{1}{\varrho}} \right)\log^{a} \frac{t}{1-t} \log^{b}(t(1-t))\D t\in\mathfrak Z_{a+b+1}(1)\tag{\ref{eq:desc3int}$'$}\label{eq:desc3int'} 
\end{align}for $ a,b\in\mathbb Z_{\geq0}$. In view of the relation\begin{align}
\left( \frac{1}{t-\varrho} +\frac{1}{t-\frac{1}{\varrho}} \right)\log^{b}(t(1-t))=b!\frac{2t-1}{1-t(1-t)}G(\boldsymbol0_{b};t(1-t))=b!\frac{\partial}{\partial t}G(1,\boldsymbol0_{b};t(1-t)),
\end{align} we may integrate by parts, and deduce the following series representation for $ \mathscr I_{a,b}$:\begin{align}
\begin{split}\mathscr I_{a,b}={}&-b!a\int_{0}^1\frac{G(1,\boldsymbol0_{b};t(1-t))}{t(1-t)}\log^{a-1} \frac{t}{1-t}\D t\\\xlongequal{\text{\cite[\S2.2]{SunZhou2026MCV}}}{}& a\int_{0}^1\sum_{k=1}^\infty\left.\!\frac{\partial^{b}}{\partial x^{b}}\frac{[t(1-t)]^{x-1}}{x}\right|_{x=k}\log^{a-1} \frac{t}{1-t}\D t\\={}& a\sum_{n=1}^\infty\int_{0}^1\left.\!\frac{\partial^{a+b-1}}{\partial \varepsilon^{a-1}\partial x^{b}}\frac{t^{x+\varepsilon}(1-t)^{x-\varepsilon}}{x}\right|_{x=n,\varepsilon=0}\frac{\D t}{t(1-t)}\\={}& a\sum_{n=1}^{\infty}\left.\!\frac{\partial^{a+b-1}}{\partial \varepsilon^{a-1}\partial x^{b}}\frac{\Gamma (x-\varepsilon) \Gamma (x+\varepsilon )}{x \Gamma (2 x)}\right|_{x=n,\varepsilon=0}.
\end{split}
\end{align} It is worth noting that when $a$ is even, all the summands vanish identically, and one has $ \mathscr I_{a,b}=0$. This is also anticipated from  the integrand of \eqref{eq:desc3int'}, which is an odd function of $2t-1$ if $a $ is even. 

In the rest of this proof, we  will focus on the series representation of  $ \mathscr I_{a,b}$ for odd  $a$, and show that it belongs to $ \mathfrak Z_{a+b+1}(1)$, using Wilf--Zeilberger (WZ) pairs \cite{WZ1990}. Thanks to Hou--Sun  (cf.\ \cite[Theorem 1.3]{HouSun2026}), we may pick a WZ pair\footnote{Here, we write $ x!\equiv\Gamma(x+1)$ for $ x>-1$. The WZ pairs listed here are just  $ \varepsilon$-shifted versions of \cite[Theorem 1.3]{HouSun2026}. }\begin{align}
\mathscr F_{\varepsilon}(n,k)\colonequals{}&\frac{2 \left(n-\frac{\varepsilon }{2}\right)! \left(k+n+\frac{\varepsilon }{2}\right)!}{3 (2 k+2 n+\varepsilon+2) (k+2 n+1)!} ,\\\mathscr G_\varepsilon(n,k)\colonequals{}& \frac{\frac{\left(n-\frac{\varepsilon }{2}+1\right)^2}{k+2 n+2}+k+n+\frac{\varepsilon }{2}+1}{n-\frac{\varepsilon }{2}+1}\mathscr F_{\varepsilon}(n,k),
\end{align}satisfying\begin{align}
\mathscr F_{\varepsilon}(n+1,k)-\mathscr F_{\varepsilon}(n,k)=\mathscr G_\varepsilon(n,k+1)-\mathscr G_\varepsilon(n,k)
\end{align}and \begin{align}
\sum_{n=0}^\infty\left.\!\frac{\partial^{a+b-1}\mathscr G_\varepsilon(x,0)}{\partial \varepsilon^{a-1}\partial x^{b}}\right|_{x=n}=\sum_{k=0}^\infty\left.\!\frac{\partial^{a+b-1}\mathscr F_\varepsilon(x,k)}{\partial \varepsilon^{a-1}\partial x^{b}}\right|_{x=0}.
\end{align}This effectively turns $ \mathscr I_{a,b}$ into a member in the $ \mathbb Q$-vector space\begin{align}
\Span_{\mathbb Q}\left\{ Z_m\sum_{k=0}^\infty\frac{1}{(k+1)^{1+s}}\prod_{j=1}^M\mathsf H_{k}^{(r_j)}\middle|\begin{smallmatrix}m+1,s,r_j\in\mathbb Z_{>0}\\Z_m\in\mathfrak Z_m(1)\\m+s+\sum_{j=1}^m r_j=a+b\\\end{smallmatrix} \right\}\subseteq\mathfrak Z_{a+b+1}(1),
\end{align}where $ \mathsf H_{k}^{(r)}\colonequals \sum_{0<n<k}\frac{1}{n^r}$ denotes the $ k$-th harmonic number of order $r$ and the ``$ \subseteq$'' step follows from \cite[Theorem 3.1(a)]{Zhou2022mkMpl}.\end{proof}
\subsection{Application to Broadhurst's conjecture}Now we verify Broadhurst's conjecture on honorary MZVs at even weights \cite[Conjecture 5]{Broadhurst2014MDV}.\begin{corollary}
For all  $w\in\mathbb Z_{>1}$, we have \begin{align}
\frac{\pi\I\Li_{w-1}(\varrho)}{3}+\R\Li_{w,1}(\varrho)\in\mathfrak Z_{w}(1).\label{eq:gw2}
\end{align}In particular, when $w$ is even, we recover  Broadhurst's relation \eqref{eq:BroadhurstConj5}.\end{corollary}\begin{proof}According to   \eqref{eq:gwn_prod_expn}, the left-hand side of  \eqref{eq:gw2} is equal to $ -\R g_{w,2}(\varrho)$, which is representable as a $ \mathbb Q$-linear combination of   MZVs at weight $w\in\mathbb Z_{>1}$.

For $ s>0$, we may express \begin{align}
\I \Li_s(\varrho)=\left( 1+\frac{1}{2^{s-1}} \right)\frac{\sqrt{3}L(\chi_{-3},s)}{2}
\end{align} through a special Dirichlet $L$-function\begin{align}
L(\chi_{-3},s)\colonequals \sum_{n=0}^\infty\left[\frac{1}{(3n+1)^s}-\frac{1}{(3n+2)^s}\right].
\end{align}If $ w$ is an even  positive  integer, then  $ L(\chi_{-3},w-1)$ is a rational multiple of $\sqrt{3}\pi^{w-1}$, which simplifies   \eqref{eq:gw2} to   $\R\Li_{w,1}(\varrho)\in\mathfrak Z_{w}(1)$.
\end{proof}
\begin{remark}At this point, one may regard Au's  \eqref{eq:Au1.5b}  as an extension of Broadhurst's depth-2 relation \eqref{eq:BroadhurstConj5} to arbitrary depths. In the meantime, we note that  Xu's relation \eqref{eq:XuConj} does not    directly transcribe the generic behavior of  $ g_{w,n}(-1)$:  there are  weights $w$ and depths $n\geq1$,  such that Au's sums $ g_{w,n}(-1)$ do not degenerate to MZVs. Symbolic computations of $\R g_{w,n}(e^{2\pi i/N})$ and $\I g_{w,n}(e^{2\pi i/N})$ at small weights (with the aid of Au's \texttt{MultipleZetaValues} package \cite{Au2022a}) also suggest the absence of Galois descents when  $N\in\{3,5,7,8,9,10,12\}$.
\eor\end{remark}


\section{Galois descents from square vertices\label{sec:desc4}}


\subsection{Multiple polylogarithms at $\frac{1+i}2$}Similar to our experience in \S\ref{subsec:desc3},  the identity\begin{align}
\begin{split}\R
g_{m+n,n}\bigg(\frac{1+i}2\bigg)={}&\left.\!\frac1{m!n!}\R\frac{\partial^{m+n}}{\partial u^{m}\partial v^n}\left[ {_2}F_1\left(\begin{array}{@{}c@{}}
1,v \\
1-u \\
\end{array}\middle| \frac{1}{i}\right) -1\right]\right|_{u=v=0}\\={}&\frac1{m!n!}\left.\!\R\frac{\partial^{m+n}}{\partial u^{m}\partial v^n}\left[ \frac{1 }{\mathrm B(1-u-v,v)} \int_0^1 \frac{t^{v-1} (1-t)^{-u-v}}{1+it} \D t-1\right]\right|_{u=v=0}
\end{split}
\end{align}reduces our task to the examination of\begin{align}
\R\frac{1}{(1+it)t}-\frac{1}{t}=-\frac{t}{1+t^2},
\end{align} a rational function whose poles are two vertices of a square inscribed in the unit circle.


\begin{theorem}We have \begin{align}
\int_0^1\frac{t}{1+t^{2}}\log^a t \log^{b+1}(1-t)\D t\in\mathfrak Z_{a+b+2}(2)\label{eq:desc4int}
\end{align}for all $a,b\in\mathbb Z_{\geq0}$.  As a result, Au's relation $\R g_{w,n}\big(\frac{1+i}2\big)\in\mathfrak Z_w(2)$ holds for all $ w\in\mathbb Z_{>0}$ and $ n\in\mathbb Z\cap[1,w]$.
\end{theorem}\begin{proof}We note that Au proposed a generalization of  \eqref{eq:desc4int}  for all $ a,b,c\in\mathbb Z_{\geq0}$ \cite[Conjecture 1.6]{Au2022a}:\begin{align}
 \int_0^1\frac{t}{1+t^{2}}\log^a t \log^{b}(1-t)\log^{c}(1+t)\D t\in\mathfrak Z_{a+b+c+1}(2),\label{eq:AuConj1.6}
\end{align}whereupon na\"ive applications of  GPL recursion \eqref{eq:GPL_rec}  and GPL shuffles \eqref{eq:GPL_shuffle} would produce summands that are in $\mathfrak Z_{a+b+c+1}(4)$. We will  prove this general conjecture of Au, by targeting its equivalent form \begin{align}\begin{split}\mathscr J_{a,b,c}\colonequals{}&\int_{0}^{1}\frac{t}{1+t^{2}}\left(\log t+\log\frac{1-t}{1+t}\right)^{a}\left(\log t-\log\frac{1-t}{1+t}\right)^{b}\log^{c}\frac{1-t^2}{t}\D t\in\mathfrak Z_{a+b+c+1}(2)
\end{split}\tag{\ref{eq:AuConj1.6}$'$}\label{eq:AuConj1.6'}
\end{align} for all $ a,b,c\in\mathbb Z_{\geq0}$.  

Without loss of generality, we momentarily assume that $ b+c$ is even and $c>0$. For any   integrable function $ f(t)$ on the open unit interval, we have \begin{align}
\int_0^1\frac{t}{1+t^{2}}\left[f(t)+f\bigg(\frac{1-t}{1+t}\bigg) \right]\D t=\int_{0}^1\frac{f(t)\D t}{1+t}\label{eq:f_01_automorph}
\end{align}by direct variable substitution. In particular, we have\begin{align}
\begin{split}&\int_{0}^{1}\frac{t}{1+t^{2}}\left(\log t+\log\frac{1-t}{1+t}\right)^{a}\left(\log t-\log\frac{1-t}{1+t}\right)^{b}\\{}&\times\left[\log^{c}\frac{1-t^2}{t}+\left( \log\frac{1-t^2}{t} -2\log2\right)^{c}\right]\D t\\={}&\int_{0}^1\frac{1}{1+t}\left(\log t+\log\frac{1-t}{1+t}\right)^{a}\left(\log t-\log\frac{1-t}{1+t}\right)^{b}\log^{c}\frac{1-t^2}{t}\D t,\label{eq:Jabc_desc_prep}
\end{split}
\end{align}where\begin{align}
\frac{1-\left(\frac{1-t}{1+t}\right)^2}{\frac{1-t}{1+t}}=\frac{4t}{1-t^2}
\end{align} and   $ \log2=-\Li_1(-1)\in\mathfrak Z_{1}(2)$.
 Applying GPL recursion \eqref{eq:GPL_rec} and GPL shuffles \eqref{eq:GPL_shuffle}  to the right-hand side of \eqref{eq:Jabc_desc_prep},  we see that it belongs to $ \mathfrak Z_{a+b+c+1}(2)$. Put differently, we have\begin{align}
\mathscr J_{a,b,c}\in \mathfrak Z_{a+b+c+1}(2)+\Span_{\mathbb Q}\left\{(\log2)^m\mathscr J_{a,b',c'}\middle|\begin{smallmatrix}m,b',c'\in\mathbb Z_{\geq0}\\c'< c\\ m+b'+c'=b+c\end{smallmatrix}\right\}\label{eq:Jabc_c_rec}
\end{align} under the working hypothesis that  $ b+c$ is even and $c>0$. Under the same  assumption, we also have the following variation on \eqref{eq:Jabc_desc_prep}:\begin{align}
\begin{split}&\int_{0}^{1}\frac{t}{1+t^{2}}\left(\log t+\log\frac{1-t}{1+t}\right)^{a}\left(\log t-\log\frac{1-t}{1+t}\right)^{b}\\{}&\times\frac{\log (1+t)\log^{c}\frac{1-t^2}{t}-\log \frac{1+t}{2}\left( \log\frac{1-t^2}{t} -2\log2\right)^{c}}{\log2}\D t\\={}&\int_{0}^1\frac{1}{1+t}\left(\log t+\log\frac{1-t}{1+t}\right)^{a}\left(\log t-\log\frac{1-t}{1+t}\right)^{b}\frac{\log (1+t)\log^{c}\frac{1-t^2}{t}}{\log2}\D t,\label{eq:Jabc_desc_prep1}
\end{split}
\end{align}where elementary algebra reveals that \begin{align}1+\frac{1-t}{1+t}={}&\frac{2}{1+t},\\
\log (1+t)={}&\log\frac{1-t^2}{t} -\left(\log t-\log\frac{1-t}{1+t}\right),
\end{align} and\begin{align}
\begin{split}&\frac{\log (1+t)\log^{c}\frac{1-t^2}{t}-\log \frac{1+t}{2}\left( \log\frac{1-t^2}{t} -2\log2\right)^{c}}{\log^{2}2}-(2c+1)\frac{\log^{c}\frac{1-t^2}{t}+\left( \log\frac{1-t^2}{t} -2\log2\right)^{c}}{2\log2}\\&{}-\frac{c}{\log2}\left( \log t-\log\frac{1-t}{1+t} \right)\left[ \log^{c-1}\frac{1-t^2}{t}+\left( \log\frac{1-t^2}{t} -2\log2\right)^{c-1} \right]\\\in{}&\mathbb Z\left[ \log2,\log\frac{1-t^2}{t},\log t-\log\frac{1-t}{1+t}\right].
\end{split}
\end{align}Here,  the left-hand side of the last equation goes like \begin{align}
c\log^{c-1}\frac{1-t^2}{t}+O\left( \log^{\max\{0,c-2\}}\frac{1-t^2}{t}\right).
\end{align}
Rewriting the right-hand side of \eqref{eq:f_01_automorph} as\begin{align}\frac12
\int_{0}^1\frac{1}{1+t}\left[f(t)+f\bigg(\frac{1-t}{1+t}\bigg) \right]\D t,
\end{align}
we have effectively shown that \eqref{eq:Jabc_c_rec} remains valid when $c$ is traded for $c-1$. In other words,  \eqref{eq:Jabc_c_rec} holds true without any constraint on the parity of    $ b+c$. This means that we may build \eqref{eq:AuConj1.6'} inductively on the situations where $c=0$.

Now we turn our attention to the claim that $ \mathscr J_{a,b,0}\in\mathfrak Z_{a+b+1}(2)$, or equivalently\begin{align}
\mathscr K_{a,b}\colonequals{}&\int_{0}^{1}\frac{t}{1+t^{2}}\log ^{a}t\log^{b}\frac{1-t}{1+t}\D t\in\mathfrak Z_{a+b+1}(2).\label{eq:Kab_defn}
\end{align}
The corresponding borderline cases are easy: we have \begin{align}
\mathscr K_{a,0}=\int_{0}^{1}\frac{t}{1+t^{2}}\log ^{a}t\D t\xlongequal{u=t^2}\frac{1}{2^{a+1}}\int_0^1\frac{\log ^{a}u\D u}{1+u}\in\mathfrak Z_{a+1}(2)
\end{align}and \begin{align}
\mathscr K_{0,b}=\int_{0}^{1}\frac{t}{1+t^{2}}\log^{b}\frac{1-t}{1+t}\D t\xlongequal{\text{\eqref{eq:f_01_automorph}}}\int_{0}^1\frac{\log ^{b}t\D t}{1+t}-\mathscr K_{b,0}\in\mathfrak Z_{b+1}(2).
\end{align}One also has the following reciprocity law: \begin{align}
\mathscr K_{a,b}+\mathscr K_{b,a}\xlongequal{\text{\eqref{eq:f_01_automorph}}}\int_{0}^1\frac{1}{1+t}\left[\log ^{a}t\log^{b}\frac{1-t}{1+t}+\log ^{b}t\log^{a}\frac{1-t}{1+t}\right]\D t\in\mathfrak Z_{a+b+1}(2).\label{eq:Krecip}
\end{align}As before, we will reduce $ \mathscr K_{a,b}$ to AMZVs and $ \mathbb Q$-linear combinations of $ Z_m\mathscr K_{a',b'}$ where $ Z_m\in\mathfrak Z_m(2)$  and $ m=a+b-a'-b'>0$. We need to distinguish several scenarios, according to the parities of $a$ and $b$, in the next   paragraph.

We evaluate an absolutely convergent integral\begin{align}
\frac{1}{(2\pi i)^{2}}\int_{i0^+-\infty}^{i0^++\infty}\frac{z}{1+z^{2}}\log ^{a+1}z\log^{b+1}\frac{z-1}{z+1}\D z
\end{align}in two ways. First, closing the contour upwards, and picking up the residue at $z=i$, we find \begin{align}
\frac{1}{4\pi i}\left( \frac{\pi i}{2} \right)^{a+b+2}.
\end{align}Second, applying a transformation $ z=1/t$ to $ \R z>1$, we get
\begin{align}
\begin{split}&\frac{1}{(2\pi i)^{2}}\int_{i0^+}^{i0^++\infty}\frac{z}{1+z^{2}}\log ^{a+1}z\log^{b+1}\frac{z-1}{z+1}\D z\\={}&\frac{1}{(2\pi i)^{2}}\int_{0}^{1}\frac{t}{1+t^{2}}\log ^{a+1}t\left[ \left(\log\frac{1-t}{1+t}+\pi i\right)^{b+1}+(-1)^{a}\log^{b+1}\frac{1-t}{1+t} \right]\D t\\{}&-\frac{(-1)^{a}}{(2\pi i)^{2}}\int_{0}^{1}\frac{1}{t}\log ^{a+1}t\log^{b+1}\frac{1-t}{1+t}\D t;\end{split}
\intertext{similarly, one has}
\begin{split}&\frac{1}{(2\pi i)^{2}}\int_{i0^+-\infty}^{i0^+}\frac{z}{1+z^{2}}\log ^{a+1}z\log^{b+1}\frac{z-1}{z+1}\D z\\={}&-\frac{(-1)^{a}}{(2\pi i)^{2}}\int_{0}^{1}\frac{t}{1+t^{2}}(\log t-\pi i)^{a+1}\left(-\log\frac{1-t}{1+t}\right)^{b+1}\D t\\{}&-\frac{1}{(2\pi i)^{2}}\int_{0}^{1}\frac{t}{1+t^{2}}(\log t+\pi i)^{a+1} \left(-\log\frac{1-t}{1+t}+\pi i\right)^{b+1}\D t\\{}&+\frac{(-1)^{a}}{(2\pi i)^{2}}\int_{0}^{1}\frac{1}{t}(\log t-\pi i)^{a+1}\left(-\log\frac{1-t}{1+t}\right)^{b+1}\D t.
\end{split}
\end{align}Reading off the real parts of the last two displayed equations, we arrive at \begin{align}
\begin{split}&\frac{[(-1)^a+1] [(-1)^b+1]}{\pi^{2}}\int_{0}^{1}\frac{ A^{a+1} B^{b+1}t\D t}{1+t^{2}}+(-1)^b (a+1)  (b+1) \int_0^1\frac{ A^{a} B^{b}t\D t}{1+t^{2}}\\&{}-\int_0^1\frac{[(-1)^a+1] (-1)^b a (a+1) B^2+ [(-1)^b+1]b (b+1)A^{2}}{2}\frac{ A^{a-1} B^{b-1}t\D t}{1+t^{2}}\\&{}-\frac{(-1)^{a}}{\pi ^{2}}\R\int_{0}^{1}\frac{A^{a+1} B^{b+1}-(A-\pi i)^{a+1}(-B)^{b+1}}{t}\D t
\\\in{}&\mathfrak Z_{a+b+1}(2)+\Span_{\mathbb Q}\left\{\pi^{2\ell}\mathscr K_{a',b'}\middle|\begin{smallmatrix}\ell\in\mathbb Z_{>0},a'\in\mathbb Z\cap[0,a],b'\in\mathbb Z\cap[0,b]\\2\ell+a'+b'=a+b\end{smallmatrix}\right\},
\end{split}
\end{align}where \begin{align}
A\colonequals\log t,\quad B\colonequals\log\frac{1-t}{1+t} .
\end{align}In particular, this implies \begin{align}
\mathscr K_{a,b}\in \mathfrak Z_{a+b+1}(2)+\Span_{\mathbb Q}\left\{\pi^{2\ell}\mathscr K_{a',b'}\middle|\begin{smallmatrix}\ell\in\mathbb Z_{>0},a'\in\mathbb Z\cap[0,a],b'\in\mathbb Z\cap[0,b]\\2\ell+a'+b'=a+b\end{smallmatrix}\right\}\label{eq:a1b1prep}
\end{align} when both  $ a$ and  $b$ are odd, as well as \begin{align}
(b+1)\mathscr K_{a,b}-a\mathscr K_{a-1,b+1}\in \mathfrak Z_{a+b+1}(2)+\Span_{\mathbb Q}\left\{\pi^{2\ell}\mathscr K_{a',b'}\middle|\begin{smallmatrix}\ell\in\mathbb Z_{>0},a'\in\mathbb Z\cap[0,a],b'\in\mathbb Z\cap[0,b]\\2\ell+a'+b'=a+b\end{smallmatrix}\right\}\label{eq:a0b1prep}
\end{align}when $a$ is even and $b$ is odd. In view of \eqref{eq:Krecip}, the last displayed equation is unscathed when its left-hand side is replaced by  $ a\mathscr K_{b+1,a-1}-(b+1)\mathscr K_{b,a}$. This effectively extends  \eqref{eq:a0b1prep} to all the cases where $ a+b$ is odd. Moreover, when  $ a+b$ is odd,  \eqref{eq:Krecip} specializes to \begin{align}
\mathscr K_{(a+b+1)/2,(a+b-1)/2}+\mathscr K_{(a+b-1)/2,(a+b+1)/2}\in\mathfrak Z_{a+b+1}(2).
\end{align}    As we combine this linearly with   \eqref{eq:a0b1prep}, we see that  \eqref{eq:a1b1prep}  is true when $ \mathscr K_{a,b}$ on its left-hand side becomes either   $ \mathscr K_{(a+b+1)/2,(a+b-1)/2}$ or $ \mathscr K_{(a+b-1)/2,(a+b+1)/2}$. After repeated invocations of  \eqref{eq:a0b1prep}, one realizes that   \eqref{eq:a1b1prep}  holds whenever the non-negative integers $ a$ and $b$ add up to an odd number.  It remains to treat the cases where both  $ a$ and  $b$ are  even. For  a fixed even number $S$ and integers $ n\in\mathbb Z\cap[0,S/2]$, consider the following analogs of  \eqref{eq:Krecip}:\begin{align}
\begin{split}&\int_{0}^1\frac{t}{1+t^{2}}\log ^{n}t\log^{n}\frac{1-t}{1+t}\left[\left(\log t+\log\frac{1-t}{1+t}\right)^{S-2n}+(-1)^{n}\left(\log t-\log\frac{1-t}{1+t}\right)^{S-2n}\right]\D t\\={}&\frac{1}{2}\int_{0}^1\frac{1}{1+t}\log ^{n}t\log^{n}\frac{1-t}{1+t}\left[\left(\log t+\log\frac{1-t}{1+t}\right)^{S-2n}+(-1)^{n}\left(\log t-\log\frac{1-t}{1+t}\right)^{S-2n}\right]\D t\\\in{}&\mathfrak Z_{S+1}(2),
\end{split}\label{eq:recip'}
\end{align} where the integrands are rational functions of $t$ times  bivariate polynomials in even powers of $\log t$ and $\log\frac{1-t}{1+t}$. When  $ S/2$ is odd, the case  $ n=S/2$ produces a trivial identity $ 0\in\mathfrak Z_{S+1}(2)$. Otherwise, one obtains linearly independent\footnote{As a quick access to the linear independence of the bivariate polynomials in question, one may map $ \log t$ to $ e^{i\phi}$ and $ \log\frac{1-t}{1+t}$ to $ e^{-i\phi}$. } integrands from \eqref{eq:recip'}, which cover   $ \mathscr K_{a,b}\in \mathfrak Z_{a+b+1}(2)$ for all the pairs of even non-negative   $ a$ and  $b$   summing to  a preset value $S$.\footnote{Replacing $ (-1)^n$ by $ (-1)^{n+1}$ in \eqref{eq:recip'}, one may also strengthen    \eqref{eq:a1b1prep} into    $ \mathscr K_{a,b}\in \mathfrak Z_{a+b+1}(2)$, when $a$ and $b$ are both odd.}
                  
Thus far, we have completed an inductive proof of \eqref{eq:Kab_defn} for all $ a,b\in\mathbb Z_{\geq0}$, which in turn, settles both   \eqref{eq:desc4int} and  \eqref{eq:AuConj1.6}. \end{proof}\subsection{Landen-type transformation}
A generalization of the Landen-type transformation formula \eqref{eq:gw1Landen} will bring us a verification of Au's \eqref{eq:Au1.5a2}.\begin{theorem}For $ w\in\mathbb Z_{>0}$ and $ n\in\mathbb Z\cap[1,w]$, we have
\begin{align}
g_{w,n}\bigg(\frac z{z-1}\bigg)={}&\sum_{j=n}^w(-1)^j\binom{j-1}{n-1}g_{w,j}(z).\label{eq:gwn_Landen}
\end{align} As a result, Au's relation $\R g_{w,n}(i)\in\mathfrak Z_w(2)$ holds for all $ w\in\mathbb Z_{>0}$ and $ n\in\mathbb Z\cap[1,w]$.
\end{theorem}\begin{proof}One may construct an equation of motion  for the left-hand side of \eqref{eq:gwn_Landen}, by applying a  variable transformation to  \eqref{eq:gwn_evolv}, before realizing that the right-hand side of \eqref{eq:gwn_Landen} evolves in the same way, with identical initial and boundary conditions.    

 Specializing  \eqref{eq:gwn_Landen} to $ z=\frac{1+i}{2}$, we get\begin{align}
\R g_{w,n}(i)=\R g_{w,n}(-i)=\sum_{j=1}^w(-1)^j\binom{j-1}{n-1}\R g_{w,j}\bigg(\frac{1+i}2\bigg),
\end{align}thereby concluding our proof. 
\end{proof}

\begin{thebibliography}{10}

\bibitem{AAR}
G.~E. Andrews, R.~Askey, and R.~Roy.
\newblock {\em Special Functions}, volume~71 of {\em Encyclopedia of
  Mathematics and Its Applications}.
\newblock Cambridge University Press, Cambridge, UK, 1999.

\bibitem{Au2022a}
K.~C. Au.
\newblock Iterated integrals and multiple polylogarithm at algebraic arguments.
\newblock \url{arXiv:2201.01676v3} [math.NT] (software available at
  \url{https://www.researchgate.net/publication/357601353_Mathematica_package_MultipleZetaValues}),
  2026.

\bibitem{BorweinStraub2015Snp}
J.~M. Borwein and A.~Straub.
\newblock Relations for {N}ielsen polylogarithms.
\newblock {\em J. Approx. Theory}, 193:74--88, 2015.

\bibitem{Broadhurst2014MDV}
D.~Broadhurst.
\newblock Multiple {D}eligne values: a data mine with empirically tamed
  denominators.
\newblock \url{arXiv:1409.7204} [hep-th], 2014.

\bibitem{Charlton2026}
S.~Charlton.
\newblock On a proof of {X}u's conjecture and depths of {G}alois descents.
\newblock {\em Rev. R. Acad. Cienc. Exactas F\'is. Nat. Ser. A Mat. RACSAM},
  120(2):Paper No. 34, 6 pages, 2026.

\bibitem{Frellesvig2016}
H.~Frellesvig, D.~Tommasini, and C.~Wever.
\newblock On the reduction of generalized polylogarithms to $\mathrm{Li}_ n$
  and $\mathrm{Li}_{2, 2}$ and on the evaluation thereof.
\newblock {\em J. High Energy Phys.}, 2016(3):189, 2016.
\newblock \url{arXiv:1601.02649v3} [hep-ph].

\bibitem{Glanois2016}
C.~Glanois.
\newblock Motivic unipotent fundamental groupoid of
  {$\mathbb{G}_m\smallsetminus\mu_N$} for {$N=2,3,4,6,8$} and {G}alois
  descents.
\newblock {\em J. Number Theory}, 160:334--384, 2016.

\bibitem{HouSun2026}
Q.-H. Hou and Z.-W. Sun.
\newblock Evaluations of some series via the WZ method.
\newblock \url{arXiv:2604.15172v1} [math.CO], 2026.

\bibitem{Koelbig1986}
K.~S. K{\"o}lbig.
\newblock Nielsen's generalized polylogarithms.
\newblock {\em SIAM J. Math. Anal.}, 17:1232--1258, 1986.

\bibitem{Panzer2015}
E.~Panzer.
\newblock Algorithms for the symbolic integration of hyperlogarithms with
  applications to {F}eynman integrals.
\newblock {\em Comput. Phys. Commun.}, 188:148--166, 2015.
\newblock (software available at
  \url{https://bitbucket.org/PanzerErik/hyperint/wiki/Home}).

\bibitem{SunZhou2026MCV}
Z.-W. Sun and Y.~Zhou.
\newblock Multiple {C}lausen values and deformed {A}p\'ery-like series.
\newblock \url{arXiv:2607.14646v2} [math.NT], 2026.

\bibitem{WZ1990}
H.~S. Wilf and D.~Zeilberger.
\newblock Rational functions certify combinatorial identities.
\newblock {\em J. Amer. Math. Soc.}, 3(1):147--158, 1990.

\bibitem{Xu2017}
C.~Xu.
\newblock Multiple zeta values and {E}uler sums.
\newblock {\em J.\ Number Theory}, 177:443--478, 2017.

\bibitem{Xu2019nonlinEuler}
C.~Xu.
\newblock Evaluations of nonlinear {E}uler sums of weight ten.
\newblock {\em Appl. Math. Comput.}, 346:594--611, 2019.

\bibitem{ZhengYang2025}
W.~Zheng and Y.~Yang.
\newblock Proof of {X}u conjecture on alternating double zeta values.
\newblock {\em Rev. R. Acad. Cienc. Exactas F\'is. Nat. Ser. A Mat. RACSAM},
  119(4):Paper No. 91, 5 pages, 2025.
\newblock (Correction: \textit{ibid.} 119(4):Paper No. 113, 2 pages, 2025).

\bibitem{Zhou2022mkMpl}
Y.~Zhou.
\newblock Hyper-{M}ahler measures via {G}oncharov--{D}eligne cyclotomy.
\newblock In H.~S. Cohl, R.~S. Costas-Santos, and R.~S. Maier, editors, {\em
  Classical hypergeometric functions and generalizations}, volume 818 of {\em
  Contemp. Math.}, pages 239--288. Amer. Math. Soc., Providence, RI, 2025.
\newblock \url{arXiv:2210.17243v4} [math.NT].

\end{thebibliography}


\end{document}
